
\subsection{Solution Design}
RD can be either stored on each level (PLRDF) or on the top level (TopLevel-RDF). When range deletes flush down to the disk, range deletes are merged and stored on to level 0 vector. Latter, when compaction happens, corresponding range deletes are moved and merged to the next level vector in PLRDF. While in TopLevel-RDF, range deletes are stored in the top level only, and those whose ranges have intersected with the newly inserted files are removed. When either flush or compaction happens, range deletes from the topper level which are newer remove the obsolete entries from the files it compacted to on the lower level. In this way, if an entry is overlapped with a range delete tombstone in a file, it must be newer than the range delete tombstone. Sequence numbers are not required because if there are overlapping files across two levels, files on the lower level is newer and if an entry and range deletes co-exist in a file, the entry is sure to exist. \\

(a) PLRDF: Range deletes are stored on each level and moved down to the next level during compaction. When a point query is blocked on a certain level, it indicates that the key doesn't exist on the levels below, requiring disk access on the current level, just the same as storing range deletes inside SST files..  \\
(b) Split-PLRDF: Whenever new entries arrive at a certain level, ranges that overlap with the newly inserted entries are split. This ensures that an entry existing on a certain level won't be covered by ranges in the RDF. However, this approach results in a larger memory footprint due to fragmented ranges. \\
(c) TopLevel-RDF: RDF is limited to the top level, such that it blocks query if an entry doesn't exist in the database. It performs split function on all incoming entries, leading to larger memory footprint. \\
